\chapter{Sequences and Families of Functions}
\label{ch:function-sequences}

For a sequence of functions, $f_k$ means the $k$th whole function,
whereas $f_k(x)$ is its scalar value at the fixed input $x$. Varying $k$
with $x$ fixed gives a sequence of numbers.

\section{Pointwise and Uniform Convergence}

For functions, convergence has two distinct quantifier patterns.  Pointwise
convergence lets the stage of approximation depend on the input; uniform
convergence does not.  This distinction explains exactly when a limit may
inherit analytic properties from its approximants.

\begin{definition}[Pointwise convergence]
Let $D\subseteq\R$, and let $f_n:D\to\R$ and $f:D\to\R$. We say that $(f_n)$ \emph{converges pointwise} to $f$ on $D$ if, for every $x\in D$,
\[
\lim_{n\to\infty}f_n(x)=f(x).
\]
\end{definition}

\begin{definition}[Uniform convergence]
\label{def:uniform-convergence}
Under the same hypotheses, $(f_n)$ \emph{converges uniformly} to $f$ on $D$ if, for every $\eps>0$, there is $N\in\N$ such that, for every $n\geq N$ and every $x\in D$,
\[
\abs{f_n(x)-f(x)}<\eps.
\]
\end{definition}

\begin{definition}[Sup norm]
If $h:D\to\R$ is bounded, define
\[
\norm{h}_\infty=\sup_{x\in D}\abs{h(x)}.
\]
Whenever $f_n-f$ is bounded, uniform convergence is equivalently
\[
\norm{f_n-f}_\infty\longrightarrow0.
\]
Indeed, the inequality $\norm{f_n-f}_\infty<\eps$ gives the same error
bound at every $x$, while a uniform error bound of $\eps/2$ forces the
supremum to be at most $\eps/2<\eps$.
\end{definition}

The dependence of the index is worth displaying:
\[
\begin{array}{ll}
\text{pointwise:}&
\forall x\in D\;\forall\eps>0\;\exists N=N(x,\eps),\\[2mm]
\text{uniform:}&
\forall\eps>0\;\exists N=N(\eps)\;\forall x\in D.
\end{array}
\]

Uniform convergence implies pointwise convergence, because one may fix $x$ after choosing the common stage $N$. The converse fails: $f_n(x)=x^n$ on $[0,1]$ converges pointwise to the discontinuous function described in Chapter~\ref{ch:motivation}, and hence cannot converge uniformly by \cref{thm:uniform-limit-continuity} below.

More directly, the pointwise limit of $x^n$ on $[0,1]$ equals $0$ on
$[0,1)$ and $1$ at $1$.  For every $n$, choosing
$x_n=\sqrt[n]{1/2}$ gives $x_n<1$ and hence $f(x_n)=0$, but
\[
\abs{f_n(x_n)-f(x_n)}=\frac12.
\]
Thus the sup-norm error never tends to zero.  In contrast,
$g_n(x)=x/n$ on $[0,1]$ converges uniformly to zero because
\[
\norm{g_n}_\infty=\frac1n.
\]

\paragraph{A working procedure.}
\[
\boxed{\text{Fix }x\text{ first; only then let }n\to\infty.}
\]
Find the pointwise candidate, separating exceptional inputs. To prove
uniform convergence, bound the error independently of $x$, or compute
$\sup_{x\in D}|f_n(x)-f(x)|$ and show it tends to zero.
To disprove it, find \emph{moving witnesses} $x_n\in D$ and a fixed
$\eps_0>0$ with $|f_n(x_n)-f(x_n)|\geq\eps_0$ for arbitrarily
large $n$. The witnesses may depend on $n$; the positive gap may not.

\begin{example}[Cases and domains]
For $f_n(x)=1/(1+nx^2)$ on $\R$, fixing $x=0$ gives the limit $1$.
For fixed $x\ne0$, $0<f_n(x)\leq1/(nx^2)\to0$.
Thus the limit has two cases. The witnesses $x_n=1/n$ have
$f(x_n)=0$ but $f_n(x_n)=1/(1+1/n)\geq1/2$, so convergence is
not uniform. On $D=\{x:|x|\geq d\}$, where $d>0$,
$\sup_D|f_n|=1/(1+nd^2)\to0$: the same formula converges uniformly
on this unbounded domain.

For $g_n(x)=x/n$, the pointwise limit is zero on $\R$.
On a bounded interval the error has a common bound:
\[
\|g_n\|_{\infty,[-A,A]}=\frac An.
\]
Choose $N>A/\eps$ for a uniform tolerance $\eps$. On $\R$, however,
the moving input $x=n$ gives $g_n(n)=1$, disproving uniform convergence. Boundedness of the domain often helps, but the preceding
example shows that an unbounded domain does not force failure.
\end{example}

\begin{example}[A uniform approximation guarantee]
Suppose a normalized input $x\in[0,1]$ should produce output
$f(x)=x^2$, while approximation level $n$ produces
$f_n(x)=x^2+x(1-x)/n$. Since
$0\leq x(1-x)=1/4-(x-1/2)^2\leq1/4$,
\[
\|f_n-f\|_\infty=\frac1{4n}.
\]
Choosing $N>1/(4\eps)$ guarantees output error below $\eps$ for
every input and every $n\geq N$. Testing only a few fixed inputs
would not establish this simultaneous guarantee.
\end{example}

\begin{proposition}[Uniform Cauchy criterion]
\label{prop:uniform-cauchy}
Let $f_n:D\to\R$. There is a function $f:D\to\R$ to which $f_n$ converges uniformly if and only if, for every $\eps>0$, there is $N$ such that, for all $m,n\geq N$ and all $x\in D$,
\[
\abs{f_n(x)-f_m(x)}<\eps.
\]
\end{proposition}
\begin{proof}
If $f_n\to f$ uniformly, choose $N$ so that
$|f_j(x)-f(x)|<\eps/2$ for every $j\geq N$ and $x\in D$.
For $m,n\geq N$ the triangle inequality gives
$|f_n(x)-f_m(x)|<\eps$ on all of $D$.
Conversely, each $(f_n(x))$ is a Cauchy real sequence, so define
$f(x)=\lim_n f_n(x)$ by completeness. Given $\eps>0$, choose $N$
from the uniform Cauchy hypothesis with tolerance $\eps/2$.
For fixed $n\geq N$ and $x\in D$, let $m\to\infty$; order
preservation gives $|f_n(x)-f(x)|\leq\eps/2<\eps$.
The same $N$ works for every $x$, proving uniform convergence.
\end{proof}

\section{Continuous Limits}

The principal benefit of uniform convergence is that one approximation index
works simultaneously near every point.  The three-term estimate below makes
this control visible and explains why pointwise convergence is insufficient.

\begin{theorem}[Uniform limit theorem]
\label{thm:uniform-limit-continuity}
Let $D\subseteq\R$. If every $f_n:D\to\R$ is continuous at $a\in D$ and $f_n\to f$ uniformly on $D$, then $f$ is continuous at $a$.
\end{theorem}
\begin{proof}
Given $\eps>0$, choose $N$ such that $\abs{f_N(x)-f(x)}<\eps/3$ for every $x\in D$. Continuity of $f_N$ at $a$ supplies $\delta>0$ such that $\abs{x-a}<\delta$ implies $\abs{f_N(x)-f_N(a)}<\eps/3$. For such $x$,
\[
\abs{f(x)-f(a)}\leq\abs{f(x)-f_N(x)}+\abs{f_N(x)-f_N(a)}+\abs{f_N(a)-f(a)}<\eps.
\]
\end{proof}

\begin{remark}[What uniform convergence permits]
Uniform limits of continuous functions are continuous by the theorem above.
Chapter~\ref{ch:riemann} proves that a uniform limit may also be passed
through a Riemann integral.  Differentiation is more delicate: uniform
convergence of $f_n$ alone does not control the difference quotients, so
Chapter~\ref{ch:differentiation} imposes uniform convergence on derivatives
and fixes one anchor value. Each interchange requires an estimate suited
to the operation: an integral uses a bound on function values, while a
derivative requires control of difference quotients.
\end{remark}

\begin{proposition}[Uniform convergence respects algebra]
\label{prop:uniform-algebra}
Let $f_n\to f$ and $g_n\to g$ uniformly on $D$. Then $f_n+g_n\to f+g$ uniformly. If both $(f_n)$ and $(g_n)$ are uniformly bounded, then
\[
f_ng_n\to fg
\]
uniformly on $D$.
\end{proposition}
\begin{proof}
For the sum choose $N_1,N_2$ for the uniform errors below $\eps/2$;
the triangle inequality works for every $n\geq\max\{N_1,N_2\}$.
For the product let
\[
\abs{f_n(x)}\leq M,
\qquad
\abs{g_n(x)}\leq K
\qquad(n\in\N,\ x\in D).
\]
Taking pointwise limits gives
\[
\abs{g(x)}\leq K
\qquad(x\in D).
\]
Then
\[
\abs{f_n(x)g_n(x)-f(x)g(x)}\leq M\abs{g_n(x)-g(x)}+K\abs{f_n(x)-f(x)}.
\]
Choose $N_1$ with $\|g_n-g\|_\infty<\eps/(2(M+1))$ for $n\geq N_1$
and $N_2$ with $\|f_n-f\|_\infty<\eps/(2(K+1))$ for $n\geq N_2$.
For every $n\geq\max\{N_1,N_2\}$ the displayed error is below
$\eps$, simultaneously at every input.
\end{proof}

\begin{proposition}[Uniform limits preserve boundedness]
\label{prop:uniform-boundedness}
Let $f_n\to f$ uniformly on $D$. If every function $f_n$ is bounded on $D$, then $f$ is bounded on $D$.
\end{proposition}
\begin{proof}
Choose $N$ so that
\[
\abs{f(x)-f_N(x)}<1
\qquad(x\in D).
\]
If $\abs{f_N(x)}\leq M$ throughout $D$, then
\[
\abs{f(x)}\leq M+1
\qquad(x\in D).
\]
\end{proof}

No common bound on the approximating functions is required: only the
chosen late function needs to supply a bound. An arbitrary early bounded
term cannot do this. On $\R$, take $f_1=0$ and $f_n(x)=x$ for $n\geq2$.
Convergence to $x$ is uniform, but the limit is unbounded.

\begin{example}[A continuous limit without uniform convergence]
On $[0,1]$, put $f_n(x)=\max\{1-n|x-1/n|,0\}$. Each graph is a continuous
triangular spike. At $0$ its value is zero; for fixed $x>0$, it is zero once
$n>2/x$. Thus $f_n\to0$ pointwise, with a continuous limit. Nevertheless
$f_n(1/n)=1$, so $\|f_n\|_\infty=1$ for every $n$. Continuity of the limit
does not supply the missing uniform control.
\end{example}

\begin{example}[Pointwise limits need not preserve boundedness]
On $D=(0,1]$, define
\[
f_n(x)=
\begin{cases}
n,&0<x\leq\dfrac1n,\\
\dfrac1x,&\dfrac1n<x\leq1.
\end{cases}
\]
Each $f_n$ is bounded by $n$, while for every $x\in D$,
\[
\lim_{n\to\infty}f_n(x)=\frac1x.
\]
The pointwise limit is unbounded. Thus the uniform hypothesis in \cref{prop:uniform-boundedness} cannot be weakened to pointwise convergence.
\end{example}

\begin{example}[Two limits need not be interchanged]
For
\[
f_n(x)=x^n
\qquad(0\leq x<1),
\]
one has
\[
\lim_{n\to\infty}\left(\lim_{x\to1^-}f_n(x)\right)=1,
\]
whereas
\[
\lim_{x\to1^-}\left(\lim_{n\to\infty}f_n(x)\right)=0.
\]
The convergence in $n$ is pointwise but not uniform near $x=1$. This is the basic obstruction behind unjustified interchanges of limits.
\end{example}

\section{Polynomial Approximation}

Polynomials are among the simplest continuous functions: they can be added,
multiplied, and differentiated exactly. How closely can they approximate
other continuous functions? Write $C([a,b])$ for the vector space of
continuous real-valued functions on $[a,b]$. Every function in this space
can be approximated on the whole interval by a polynomial, to any
prescribed uniform tolerance.

\subsection*{Reading the Bernstein weights}
For degree two the approximation is
\[
B_2(f)(x)=f(0)(1-x)^2+2f\left(\frac12\right)x(1-x)+f(1)x^2.
\]
At $x=1/2$ its weights are $1/4,1/2,1/4$. For $f(t)=t^2$ this
gives $B_2(f)(x)=(x+x^2)/2$, not $x^2$ itself. The approximating
polynomial samples nearby values heavily; exact reproduction of every
quadratic is not the aim. The proof below establishes concentration using
finite algebraic sums, with no probability prerequisite.
\begin{theorem}[Weierstrass approximation theorem]
\label{thm:weierstrass-approximation}
If $f\in C([a,b])$, then for every $\eps>0$ there is a polynomial $p$ such
that
\[
\norm{f-p}_\infty<\eps.
\]
Equivalently, polynomials are uniformly dense in $C([a,b])$.
\end{theorem}
\begin{proof}
We first work on $[0,1]$.  For $f\in C([0,1])$, define its $n$th Bernstein
polynomial by
\[
B_n(f)(x)=\sum_{k=0}^nf\left(\frac kn\right)\binom nkx^k(1-x)^{n-k}.
\]
It is a polynomial.  Its nonnegative weights have sum $1$ by the binomial
theorem, so $B_n(f)(x)$ is a weighted average of sampled values $f(k/n)$.
Before estimating a general function, check the samples $1,t,t^2$:
\[
B_n(1)=1,\qquad B_n(t)=x,\qquad B_n(t^2)=x^2+\frac{x(1-x)}n.
\]
For weights $w_{n,k}=\binom nkx^k(1-x)^{n-k}$, the first moment is
\[
\sum_{k=0}^n\frac{k}{n}w_{n,k}
=x\sum_{k=1}^n\binom{n-1}{k-1}x^{k-1}(1-x)^{n-k}=x.
\]
The factorial second moment is similarly
$\sum k(k-1)w_{n,k}=n(n-1)x^2$. Since $k^2=k(k-1)+k$,
division by $n^2$ gives the displayed quadratic identity.

These follow without differentiation: use
$k\binom nk=n\binom{n-1}{k-1}$ and
$k(k-1)\binom nk=n(n-1)\binom{n-2}{k-2}$, sum with the binomial theorem
from \cref{thm:appendix-binomial},
and write $k^2=k(k-1)+k$ (the case $n=1$ is direct).
Expanding the square now shows that the weights concentrate near $k/n=x$:
\[
\sum_{k=0}^n\left(\frac{k}{n}-x\right)^2
\binom nkx^k(1-x)^{n-k}=\frac{x(1-x)}n\leq\frac1{4n}.
\]

Let $M=\norm f_\infty$.  Given $\eps>0$, uniform continuity of $f$ supplies
$\delta>0$ such that $|s-t|<\delta$ implies
$|f(s)-f(t)|<\eps/2$.  Split the defining weighted average for
$B_n(f)(x)-f(x)$ according as $|k/n-x|<\delta$ or not.  The first part has
absolute value at most $\eps/2$.  On the second part, the error is at most
$2M$, while
\[
\sum_{|k/n-x|\geq\delta}\binom nkx^k(1-x)^{n-k}
\leq\frac1{\delta^2}\sum_{k=0}^n\left(\frac{k}{n}-x\right)^2
\binom nkx^k(1-x)^{n-k}
\leq\frac1{4n\delta^2}.
\]
Choose $N>\max\{1,M/(\eps\delta^2)\}$. For every $n\geq N$,
the second contribution is at most $M/(2n\delta^2)<\eps/2$,
uniformly in $x$.  Hence $B_n(f)\to f$ uniformly on $[0,1]$.

For a general interval $[a,b]$, apply this result to
$F(t)=f(a+(b-a)t)$ and substitute $t=(x-a)/(b-a)$.  The resulting expression
is a polynomial in $x$ and has the required uniform error.
\end{proof}

The error estimate separates two reasons a contribution can be small:
\[
\begin{array}{c|c|c}
\text{samples}&\text{function error}&\text{total weight}\\\hline
|k/n-x|<\delta&<\eps/2&\leq1\\
|k/n-x|\geq\delta&\leq2M&\leq1/(4n\delta^2)
\end{array}
\]
The near samples are accurate because of uniform continuity. The far
samples need not be accurate, but together carry little weight because
each contributes at least $\delta^2$ times its weight to the squared
distance sum. First choose $\delta$ from $f$, then choose $N$ as above.
Neither choice depends on the evaluation point $x$.
\begin{example}[Approximating a corner]
The function $f(x)=\abs{x}$ on $[-1,1]$ is not a polynomial and is not
differentiable at $0$.  Nevertheless, the theorem provides polynomials
$p_n$ with $\norm{p_n-\abs{x}}_\infty\to0$.  This does not contradict the
fact that uniform limits of differentiable functions need not be
differentiable: uniform convergence preserves continuity, but derivatives
require the stronger hypotheses discussed in Chapter~\ref{ch:differentiation}.
\end{example}

\begin{remark}[Optional perspective: Stone--Weierstrass]
For a compact Hausdorff space $X$, the real Stone--Weierstrass theorem says
that a subalgebra of $C(X,\R)$ which contains the constants and separates
points is dense in the sup norm.  On $[a,b]$, polynomials have these
properties, so the classical theorem is its first example.  The general
result belongs naturally to functional analysis and is not used here.
\end{remark}

\section{Equicontinuity}

Uniform convergence controls variation in the sequence index.  Equicontinuity
instead controls variation in the input uniformly across a family; together
with compactness it is a central compactness principle for function spaces.

\begin{definition}[Equicontinuity]
Let $\mathcal F$ be a family of functions $D\to\R$. It is \emph{equicontinuous at $a\in D$} if, for every $\eps>0$, there is $\delta>0$ such that every $f\in\mathcal F$ satisfies
\[
\abs{x-a}<\delta\Longrightarrow\abs{f(x)-f(a)}<\eps
\qquad(x\in D).
\]
It is equicontinuous on $D$ if it is equicontinuous at every point of $D$.
\end{definition}

\begin{proposition}
If $f_n\to f$ pointwise on $D$ and the family $\set{f_n:n\in\N}$ is equicontinuous at $a$, then $f$ is continuous at $a$.
\end{proposition}
\begin{proof}
Choose the common input neighborhood first, then take
pointwise limits at the two fixed points. Equicontinuity gives $\delta$
such that $|f_n(x)-f_n(a)|<\eps/2$ for every $n$ whenever
$x\in D$ and $|x-a|<\delta$. Fix such $x$ and let $n\to\infty$.
The inequality becomes $|f(x)-f(a)|\leq\eps/2<\eps$.
This is continuity at $a$; no uniform convergence was needed.
\end{proof}

\begin{example}[A shared Lipschitz bound and moving spikes]
If $|f(x)-f(y)|\leq L|x-y|$ for every member of a family, then
$\delta=\eps/(L+1)$ proves equicontinuity at every point.
The functions $x\mapsto c+sx$ with $|s|\leq3$ give an example even
when their constants $c$ are unbounded. Equicontinuity controls variation,
not absolute heights. The triangular spikes introduced earlier are not
equicontinuous at zero: $f_n(0)=0$, $f_n(1/n)=1$, and $1/n\to0$.
Every proposed common $\delta$ fails for a late spike.
\end{example}
\begin{lemma}[Uniform equicontinuity on a compact interval]
\label{lem:uniform-equicontinuity-interval}
Let $\mathcal F$ be an equicontinuous family of functions
\[
[a,b]\longrightarrow\R.
\]
Then, for every $\eps>0$, there is $\delta>0$ such that
\[
\abs{x-y}<\delta
\quad\Longrightarrow\quad
\abs{f(x)-f(y)}<\eps
\]
for every $f\in\mathcal F$ and every $x,y\in[a,b]$.
\end{lemma}
\begin{proof}
If not, there would be $\eps>0$, functions $h_n\in\mathcal F$, and points
$x_n,y_n\in[a,b]$ such that
\[
\abs{x_n-y_n}<\frac1n,
\qquad
\abs{h_n(x_n)-h_n(y_n)}\geq\eps.
\]
By compactness of $[a,b]$, pass to a subsequence with $x_n\to x$.  Then also
$y_n\to x$.  Relabel this subsequence as $(x_n,y_n,h_n)$.  Equicontinuity at
$x$, with tolerance $\eps/3$, supplies $\delta>0$ working for every $h_n$.
Choose $N_1,N_2$ so that $|x_n-x|<\delta$ for $n\geq N_1$ and
$|y_n-x|<\delta$ for $n\geq N_2$. For $n\geq\max\{N_1,N_2\}$,
both $|h_n(x_n)-h_n(x)|$ and $|h_n(y_n)-h_n(x)|$ are below $\eps/3$.  The triangle inequality gives a contradiction.
\end{proof}

\begin{theorem}[Arzel\`a--Ascoli on an interval]
\label{thm:arzela-ascoli-interval}
Let $(f_n)$ be an equicontinuous sequence in $C([a,b])$ for which there is
$M>0$ satisfying
\[
\abs{f_n(x)}\leq M
\qquad(n\in\N,\ x\in[a,b]).
\]
Then it has a subsequence that converges uniformly on $[a,b]$ to a
continuous function.
\end{theorem}
\begin{proof}
The idea is to first obtain convergence on a countable dense set and then
use equicontinuity to make that convergence uniform.  Enumerate a countable
dense subset of $[a,b]$ as $(q_j)$.  Uniform boundedness and
Bolzano--Weierstrass allow us to choose successively a subsequence converging
at $q_1$, then a subsequence of it converging at $q_2$, and so on.  The
diagonal subsequence can be made explicit. Write the nested increasing
index lists as $(n^{(j)}_k)_{k\geq1}$ and take $g_k=f_{n^{(k)}_k}$.
Because list $k+1$ is a subsequence of list $k$,
$n^{(k+1)}_{k+1}\geq n^{(k)}_{k+1}>n^{(k)}_k$. For fixed $j$, all entries
with $k\geq j$ belong to list $j$, in increasing order. Thus $(g_k(q_j))$
converges. The remaining step upgrades these separate convergences to one
uniform estimate by using only finitely many nearby test points.

The diagonal selection can be pictured as follows; each row is a
subsequence of the preceding row, and the marked diagonal gives $g_k$:
\[
\begin{array}{c|cccc|l}
&1&2&3&\cdots&\text{convergence guaranteed at}\\\hline
1&\boxed{n^{(1)}_1}&n^{(1)}_2&n^{(1)}_3&\cdots&q_1\\
2&n^{(2)}_1&\boxed{n^{(2)}_2}&n^{(2)}_3&\cdots&q_1,q_2\\
3&n^{(3)}_1&n^{(3)}_2&\boxed{n^{(3)}_3}&\cdots&q_1,q_2,q_3
\end{array}
\]

Given $\eps>0$, \cref{lem:uniform-equicontinuity-interval} supplies
$\delta>0$ such that
\[
\abs{x-y}<\delta
\quad\Longrightarrow\quad
\abs{g_k(x)-g_k(y)}<\eps/3
\qquad(k\in\N).
\]
Using density and a finite partition of $[a,b]$, choose finitely many of the
$q_j$ so that every $x\in[a,b]$ lies within $\delta$ of one of them.  List the selected points as $q_{j_1},\ldots,q_{j_\ell}$.
At $q_{j_i}$ choose a Cauchy stage $N_i$ for tolerance $\eps/3$.
Put $N=\max\{N_1,\ldots,N_\ell\}$. For $r,s\geq N$ the values
at every selected point differ by less than $\eps/3$.  For a nearby selected point $q$, the
three-term estimate is
\[
|g_r(x)-g_s(x)|\leq
\underbrace{|g_r(x)-g_r(q)|}_{\text{equicontinuity}}
+\underbrace{|g_r(q)-g_s(q)|}_{\text{finite test points}}
+\underbrace{|g_s(q)-g_s(x)|}_{\text{equicontinuity}}.
\]
There are only finitely many test points, so taking the maximum of their
Cauchy stages gives one $N$. The triangle inequality then gives
\[
\abs{g_r(x)-g_s(x)}<\eps
\qquad(x\in[a,b]).
\]
Thus $(g_k)$ is uniformly Cauchy, so \cref{prop:uniform-cauchy} gives a
uniform limit.  It is continuous by \cref{thm:uniform-limit-continuity}.
\end{proof}



\section*{Exercises}
\addcontentsline{toc}{section}{Exercises}

\begin{hexercise}{uniform-cases}
Find the pointwise limit of $x^n$ on $[-1,1]$ wherever it exists. Determine uniform convergence on $[0,r]$ for $0<r<1$, on $[0,1)$, and on $[0,1]$.
\end{hexercise}

\begin{exercise}
For $f_n(x)=x/(1+nx)$ on $[0,\infty)$, find the pointwise limit and an input-independent error bound. Does unboundedness of the domain prevent uniform convergence?
\end{exercise}

\begin{hexercise}{uniform-model}
On a normalized input interval $[0,1]$, an approximation adds $x(1-x)/n$ to the desired output. Derive a worst-case error, and choose a level guaranteeing error below $10^{-3}$ for every input.
\end{hexercise}

\begin{exercise}
Prove directly that a uniform limit of functions with a shared Lipschitz constant $K$ has the same Lipschitz constant. Can a uniformly convergent sequence have Lipschitz constants tending to infinity?
\end{exercise}

\begin{exercise}
On $[0,1]$, let $f_n(x)=x^n$.
Determine its pointwise limit and prove directly from \cref{def:uniform-convergence} that the convergence is not uniform.
\end{exercise}

\begin{exercise}
Let $f_n(x)=x/n$ for $x\in\R$.
Prove that $f_n$ converges pointwise to zero but not uniformly on $\R$.
\end{exercise}

\begin{exercise}
Suppose that $f_n\to f$ uniformly on $D$ and that every $f_n$ is bounded by the same number $M$. Prove that
\[
\abs{f(x)}\leq M
\qquad(x\in D).
\]
\end{exercise}

\begin{hexercise}{moving-spike}
For $f_n(x)=nx/(1+n^2x^2)$ on $[0,1]$, find the pointwise limit and
a moving point preventing uniform convergence. Prove uniform convergence
on $[\delta,1]$ for each $\delta>0$.
\end{hexercise}
\begin{exercise}
Compute $B_2(f)$ for $f(t)=|2t-1|$. Compare its error at $0,1/2,1$.
Which part of the Bernstein argument depends on continuity, and which
part only on the binomial identities?
\end{exercise}
\begin{exercise}
Show that pointwise convergence to $f$ and equicontinuity on $[a,b]$
imply uniform convergence there. Use the continuity of $f$ proved above
and a finite set of nearby test points. Contrast the moving-spike example.
\end{exercise}
\begin{exercise}
Give a uniformly bounded family that is not equicontinuous, and an
equicontinuous family that is not uniformly bounded. Explain why neither
alone supplies the hypotheses of Arzel\`a--Ascoli.
\end{exercise}
\begin{exercise}
Suppose $f_n(0)=0$ and $|f_n(x)-f_n(y)|\leq2|x-y|$ on $[0,1]$.
Use Arzel\`a--Ascoli and prove that every uniform subsequential limit
also has this Lipschitz bound. Must the original sequence converge?
\end{exercise}
